\centerline{\bf \Huge Степень вхождения}

\begin{flushright}
        \scriptsize{Эпизод 5. Рей, по ту сторону сердца.}
\end{flushright}

\problem{1}
Даны натуральные числа $m, n, k$. Оказалось, что $m^k+n^k+m$ делится на $m n$. Докажите, что $m$ - точная $k$-я степень натурального числа.

\problem{2}
Натуральные числа $b$ и $n$, большие 1 , таковы, что для каждого натурального $k>1$ существует натуральное число $a_k$ такое, что $b-a_k^n$ делится на $k$. Докажите, что $b=A^n$ для некоторого натурального $A$.

\problem{3}
Найдите все натуральные $m$ такие, что $3^3+6^6+\ldots+(3 m)^{3 m}$ является квадратом натурального числа.

\problem{4}
Для натурального числа $n$ обозначим через $S_n$ наименьшее общее кратное всех чисел $1,2,3, \ldots, n$. Существует ли такое натуральное число $m$, что $S_{m+1}=4 S_m$ ?

\problem{5}
Петя выбрал несколько последовательных натуральных чисел и каждое записал либо красным, либо синим карандашом (оба цвета присутствуют). Может ли сумма наименьшего общего кратного всех красных чисел и наименьшего общего кратного всех синих чисел являться степенью двойки?

\problem{6}
Для натурального числа $n$ есть ровно два из чисел $1,2, \ldots, 2022$, на которые оно не делится. Пусть эти числа - $a$ и $b$. Может ли оказаться, что $a-b=13$ ?

\problem{7}
Натуральные числа $a, b, c$ таковы, что $\frac{a}{b}+\frac{b}{c}+\frac{c}{a}$ - целое. Докажите, что $a b c-$ точный куб.

\problem{8}
Тройка целых чисел $(x, y, z)$, наибольший общий делитель которых равен 1, является решением уравнения

$$
y^2 z+y z^2=x^3+x^2 z-2 x z^2 .
$$


Докажите, что $z$ является кубом целого числа.


\newpage

\hspace{3mm}

\problem{9}
Пусть $p>2$ - натуральное число, не делящееся на 3. Рассмотрим все целые числа $a_1, a_2, \ldots, a_k$ из интервала $\left(-\frac{p}{2}, \frac{p}{2}\right)$ такие, что $a_i \equiv p\pmod{3}$ для всех $i=1,2,3, \ldots, k$. Докажите, что произведение

$$
\frac{p-a_1}{\left|a_1\right|} \cdot \frac{p-a_2}{\left|a_2\right|} \cdot \ldots \cdot \frac{p-a_k}{\left|a_k\right|}
$$

является натуральной степенью тройки.

\problem{10}
\begin{enumerate}
    \item[a)] Даны 10 натуральных чисел $a_1, a_2, \ldots, a_{10}$. Известно, что $a_1+a_2+\ldots+a_{10}=1000$. Оказалось, что произведение их факториалов
        
    $$
    a_{1}!\cdot a_{2}!\cdot \ldots \cdot a_{10}!
    $$
    
    
    является 10-й степенью натурального числа. Докажите, что вседанные числа равны.
    \item[b)] 
    Let $a_1, a_2, \ldots , a_k$ positive integers such that $a_1+a_2+ \ldots + a_k = k(p-1),$ where $p$ is prime number and $k\ > 2$, then the product $a_1! \cdot a_2! \cdot \ldots \cdot a_k!$ will be an exact $k$-th power of a natural number if and only if 
    \[a_1=a_2=\ldots=a_k=p-1.\]
\end{enumerate}

В системе счисления по основанию 2026 записано число

$$
N=\overline{1\ 2\ 3 \ldots 2025}_{2026}
$$

(то есть его цифры подряд: $1,2,3, \ldots, 2025$ ).
Найдите две последние цифры числа $N$ в системе счисления по основанию 45.


В системе счисления по основанию 45 записано число

$$
N=\underbrace{20262026 \ldots 2026}_{2026 \text { раз подряд }} 45
$$

(то есть блок из двух цифр 20, 26 повторён ровно 2026 раз).
Найдите последнюю цифру числа $N$ в системе счисления по основанию 2026.